\chapter{结论与展望}

\section{结论}
写到这里，这篇毕业论文也算接近尾声了。回顾整个项目开发过程，本课题的核心目标是解决高校经费管理中决策手段匮乏、预警不及时的问题。通过设计并实现这套决策支持系统，笔者认为核心成果主要体现在以下几个方面：

（1）构建了完整的全栈技术架构：基于 Spring Boot 2.6.13 与 MyBatis-Plus 实现了稳健的后端业务链路，配合 Vue 3 与 ECharts 构建了交互式决策支持看板。通过数据采集、持久化到多维可视化的闭环，解决了科研经费数据“孤岛化”的问题。

（2）实现了精准的风险预警机制：系统通过“经费执行率”与“自然时间进度”的二元对比模型，结合智能标签引擎，实现了对执行严重滞后、预算余额告急以及逾期未结等风险点的实时扫描与软预警。

（3）强化了经费管理的业务闭环：在经费报销模块中引入了科目余额硬拦截算法，确保每一笔报销均在受控范围内。系统不仅满足了基础的账目维护需求，更显著提升了管理决策的科学性与透明度。

\section{展望}
虽然系统目前已经能够满足核心业务需求并提供初步的决策支持，但在实际应用的广度和深度上仍存在不少改进空间。

移动化办公是一个重要的发展方向。目前的科研人员普遍处于快节奏的工作状态，如果能开发配套的移动端应用，让他们在手机上就能实时监控项目进度并接收风险提醒，将进一步提升系统的实用性。

目前的预警机制主要基于固定的业务阈值，灵活性仍有欠缺。在未来的研究中，可以考虑引入机器学习算法，对海量历史支出数据进行深度挖掘，从而实现更精准的支出趋势预测和异常诊断。

最后，系统目前的数据来源相对单一。未来若能与高校现有的资产管理系统、采购系统实现深层对接，构建起跨部门的数据中台，那么决策分析的维度和准确度将会得到质的飞跃。
